\chapter*{Cómo usar este manual}
\addcontentsline{toc}{chapter}{Cómo usar este manual}
\markboth{Cómo usar este manual}{Cómo usar este manual}

Este manual es para cualquier persona que quiera usar Lázaro Linux, aunque
nunca haya usado otra cosa que Windows o el celular. No hace falta saber
de computación: cada palabra técnica se explica la primera vez que
aparece, y todas están juntas en el \hyperref[cap:a-glosario]{Glosario} del
final.

\paragraph{No hace falta leerlo entero.} Está dividido en cuatro partes:
\begin{itemize}
  \item \textbf{Empezar} (capítulos 1 a 5): qué es Lázaro, cómo se prepara el
        USB y cómo se instala. Léela de corrido la primera vez.
  \item \textbf{Uso diario} (capítulos 6 a 14): internet, documentos, PDF,
        videos, impresoras, programas. Ve directo a lo que necesites.
  \item \textbf{Herramientas de Lázaro} (capítulo 15): los programas propios
        de Lázaro para que el equipo vaya mejor.
  \item \textbf{Cuando algo falla} (capítulos 16 a 18): qué hacer si algo
        no funciona y cómo pedir ayuda.
\end{itemize}

\paragraph{Cómo se escriben las cosas.}
\begin{itemize}
  \item Los botones, menús y opciones que ves en la pantalla van así:
        \menu{Archivo > Guardar como}. Quiere decir: abre el menú
        \emph{Archivo} y elige \emph{Guardar como}.
  \item Las teclas van así: \keys{\ctrl + C} quiere decir \enquote{mantén
        apretada la tecla Ctrl y, sin soltarla, pulsa la C}.
  \item Lo que se escribe en la terminal va en un recuadro oscuro. No hace
        falta usar la terminal para casi nada, pero a veces es el camino más
        corto:
\end{itemize}
\begin{terminal}
lazaro-afinador --revisar
\end{terminal}

\paragraph{Los recuadros de colores.}
\begin{consejo}
Un truco que te ahorra tiempo.
\end{consejo}
\begin{cuidado}
Algo que puede hacerte perder archivos o dejar el equipo sin arrancar.
Léelo siempre.
\end{cuidado}
\begin{nota}
Un dato extra para quien quiera entender más. Puedes saltarlo.
\end{nota}

\paragraph{Las fuentes.} Al final de cada capítulo hay una lista de páginas
de confianza (la documentación oficial de Alpine, de cada escritorio y de
cada programa) para profundizar. Los números entre corchetes en el texto,
como \cite{alpine-about}, llevan a esa lista. Todas juntas están en la
Bibliografía, al final del libro.

\paragraph{Escritorios distintos.} Lázaro ofrece varios escritorios (el
\enquote{aspecto} del sistema). Las capturas de pantalla son de
\textbf{XFCE}, el que viene marcado de fábrica. Si elegiste otro, los
nombres de los menús pueden cambiar un poco; el capítulo~\ref{cap:06-escritorio}
explica las diferencias.
